%% ma: L11-B17-0009
%% lop: 11
%% chuong: 5
%% bai: 17
%% dang: 11.17.2
%% loai: TL
%% muc: TH
%% dap_an: "-"
%% nguon: "tai-lieu-hoc-tap-mon-toan-lop-11-theo-chuong-trinh-giao-duc-pho-thong-2018.pdf (Lớp Toán Thầy Duy, Phạm Lê Duy, THPT Chu Văn An), Chương 3 Bài 3, Dạng 2, Ví dụ 9 (tr.390-392)"
%% kien_thuc: "Lớp 11 Bài 17 (hàm số liên tục)"
%% trang_thai: cho-duyet
%% ly_do: "Dạng mới đề xuất (11.17.x), chờ giáo viên duyệt dạng; tài liệu gốc không có lời giải mẫu, đáp án do kho tự giải và kiểm bằng sympy."
\begin{ex}
Cho hàm số $f(x)=\begin{cases}(x+1)^2&\text{khi }x>1\\x^2+3&\text{khi }x<1\\k^2&\text{khi }x=1\end{cases}$. Tìm $k$ để $f(x)$ gián đoạn tại $x=1$.
\loigiai{\sloppy $\lim\limits_{x\to1^+}f(x)=4$ và $\lim\limits_{x\to1^-}f(x)=4$ nên $\lim\limits_{x\to1}f(x)=4$. Hàm số gián đoạn tại $x=1$ khi $f(1)=k^2\ne4\Leftrightarrow k\ne\pm2$.}
\end{ex}
